% !TeX root = ../../ONMIPA Alin.tex
Struktur ruang vektor memungkinkan kita membentuk kombinasi linear,
memilih basis, dan menghitung dimensi. Hasil kali dalam menambahkan
gagasan geometris: panjang, jarak, dan ortogonalitas. Dengan struktur
ini, kita dapat menentukan vektor terdekat dalam suatu subruang
serta memilih basis yang menyederhanakan perhitungan.

Pembahasan dimulai dari aksioma dan ketaksamaan dasar, dilanjutkan
dengan basis ortonormal, proyeksi, serta pendekatan kuadrat terkecil.
Pada bagian akhir, hasil kali dalam digunakan untuk mempelajari
adjoin, isometri, dan operator adjoin diri. Pembaca diharapkan telah
memahami basis, dimensi, kernel, citra, dan representasi matriks
dari transformasi linear.
